% ===================================================================
%  圖表第 15 號 — 街市。--- 食物。
%
%  Traduction, faite en 2026, en cantonais ecrit d'aujourd'hui, en
%  caracteres traditionnels, sur l'ido de texto/io. Regles de la
%  colonne : voir 10-toubiu-01.tex. Une seule numerotation.
%
%  C'EST LE TABLEAU OU LES DEUX LANGUES SE SEPARENT LE PLUS, ET LE
%  PLUS SIMPLEMENT : sur ce qui se mange. Pas une syntaxe, pas un
%  emprunt — des noms differents pour les memes legumes.
%
%    番茄 (14)   la tomate    Pekin : 西红柿
%    椰菜花 (15) le chou-fleur Pekin : 菜花
%    椰菜 (17)   le chou      Pekin : 卷心菜
%    翠玉瓜 (19) la courgette Pekin : 西葫芦
%    薯仔 (61)   la pomme de terre Pekin : 土豆
%    甘筍        la carotte   Pekin : 胡萝卜
%    芫茜        le persil    Pekin : 香菜
%    西洋菜      le cresson
%    提子 (54)   le raisin    Pekin : 葡萄
%    合桃 (57)   la noix      Pekin : 核桃
%    青口        la moule     Pekin : 贻贝
%    吞拿魚      le thon      Pekin : 金枪鱼
%    沙甸魚 (50) la sardine   Pekin : 沙丁鱼
%    蠔 (98)     l'huitre     Pekin : 牡蛎
%    牛扒        le bifteck   Pekin : 牛排
%    朱古力 (34) le chocolat  Pekin : 巧克力
%    番梘 (38)   le savon     — deja releve au tableau 6
%
%  ET LE MARCHE LUI-MEME EST 街市, « le marche de rue », la ou Pekin
%  dit 菜市场, « la place aux legumes ». Le mot de Hong Kong nomme le
%  lieu ; celui du nord nomme la marchandise.
%
%  DEUX MOTS DE CE TABLEAU N'ONT PAS ETE CHOISIS, ILS SE SONT
%  IMPOSES. Le bloc 16 raconte l'agent (100) qui pourchasse les
%  vendeuses sans patente (99) : le cantonais de Hong Kong a un mot
%  pour exactement cela, 走鬼 — « courir au fantome », ce que crie le
%  colporteur quand la police arrive, et par extension le colporteur
%  lui-meme. L'ido dit « legum-kolportistini », le francais « petites
%  marchandes » ; le cantonais dit ce que la scene est. Et les
%  clientes du bloc 4 sont 師奶, le mot ordinaire de la maitresse de
%  maison qui fait son marche.
%
%  LE TABLEAU N'A DEMANDE AUCUNE INVERSION. Il enumere du debut a la
%  fin, et une enumeration suit l'ordre qu'on lui donne — c'est le
%  premier de la colonne a sortir juste sans une seule reprise.
% ===================================================================

%%K t15-apar-1 apar
\VUsaut{4.0mm}
\VUtitre{15.0pt}{60}{圖表第 15 號}
\VUsaut{3.0mm}

%%K t15-tit-1 sub
\VUpk{11.4pt}{40}{街市。--- 食物。}
\VUsaut{3.0mm}

%%K t15-01-1 p
1. --- 供應我哋食用嘅商家，大部分都聚喺\VUgras{有蓋街市}個\VUgras{大棚}
\textsuperscript{(1)} 底下，或者喺隔籬幾條街度。

%%K t15-02-1 p
2. --- \VUgras{豬肉佬} \textsuperscript{(2)} 賣\VUgras{火腿} \textsuperscript{(3)}、
\VUgras{血腸} \textsuperscript{(4)}、\VUgras{香腸} \textsuperscript{(5)}、\VUgras{生腸}
\textsuperscript{(6)} 同\VUgras{肥豬肉} \textsuperscript{(7)}；佢個檔就喺
\VUgras{賣牛油芝士嘅婆婆} \textsuperscript{(8)} 隔籬。婆婆個檔擺住紅皮嘅
\VUgras{荷蘭芝士} \textsuperscript{(9)}、\VUgras{卡蒙貝爾芝士} \textsuperscript{(10)}、
\VUgras{格魯耶爾芝士餅} \textsuperscript{(11)}，同新鮮嘅\VUgras{牛油磚} \textsuperscript{(12)}。

%%K t15-03-1 p
3. --- 婆婆對面，\VUgras{賣菜婆} \textsuperscript{(13)} 向啲女客推銷靚\VUgras{番茄}
\textsuperscript{(14)}、\VUgras{椰菜花} \textsuperscript{(15)}、\VUgras{生菜} \textsuperscript{(16)}、
\VUgras{椰菜} \textsuperscript{(17)}、\VUgras{翠玉瓜} \textsuperscript{(19)}、\VUgras{蜜瓜}
\textsuperscript{(18)}，連\VUgras{蘑菇} \textsuperscript{(20)} 都有。不過有個
\VUgras{送啤酒嘅}將架\VUgras{送貨車} \textsuperscript{(21)} 停正喺佢檔口前面，車上載滿
\VUgras{啤酒桶} \textsuperscript{(22)}，佢啲客就埋唔到嚟。有個種菜婆捧住一細籮
\VUgras{洋薊} \textsuperscript{(23)} 送過嚟畀佢。

%%K t15-tit-2 sub
\VUsaut{4.0mm}
\VUpk{10.2pt}{40}{買同賣。}
\VUsaut{3.0mm}

%%K t15-04-1 p
4. --- 街市入面熱鬧到不得了，因為夠鐘\VUgras{拍賣} \textsuperscript{(24)} 嘞。啲
\VUgras{師奶} \textsuperscript{(26)} 嚟買餸，行過\VUgras{賣豬雜婆} \textsuperscript{(25)} 個檔口
就停低，買啲\VUgras{牛柏葉}或者\VUgras{牛仔頭。}

%%K t15-05-1 p
5. --- 有個肥肥哋嘅\VUgras{大廚} \textsuperscript{(27)} 喺\VUgras{賣魚婆}
\textsuperscript{(28)} 檔口度講緊一條靚\VUgras{吞拿魚}嘅價；個賣魚婆就話幾多就幾多，
一路抬價。佢今日收到好多\VUgras{鱔、撻沙、鱒魚}同\VUgras{鯉魚。}佢啲\VUgras{鱈魚}
同\VUgras{青口}都好新鮮。

%%K t15-tit-3 sub
\VUsaut{4.0mm}
\VUpk{10.2pt}{40}{平嘢同貴嘢。}
\VUsaut{3.0mm}

%%K t15-06-1 p
6. --- 佢隔籬個\VUgras{賣雞鴨婆} \textsuperscript{(29)} 就好有良心。佢賣\VUgras{火雞、}
\VUgras{鵝、鴨}或者普通\VUgras{雞項，}都係照實價收。

%%K t15-07-1 p
7. --- 廣場個角落頭，一樓，最近開咗個\VUgras{賣布佬} \textsuperscript{(30)}；啲女客
喺佢度實搵到好靚嘅\VUgras{檯布、餐巾、圍裙}同\VUgras{抹布。}同一層樓，
有個\VUgras{磨刀佬} \textsuperscript{(31)} 開咗檔。佢幫街市啲檔主磨\VUgras{刀，}又用最硬嘅
\VUgras{鋼}幫啲種菜佬打\VUgras{鐮刀。}

%%K t15-tit-4 sub
\VUsaut{4.0mm}
\VUpk{10.2pt}{40}{雜貨。}
\VUsaut{3.0mm}

%%K t15-08-1 p
8. --- 佢個檔下面，最近開咗間好大嘅食品舖。佢已經唔係舊時嗰啲細細間、
唔起眼嘅\VUgras{雜貨舖} \textsuperscript{(32)} 嘞——嗰啲淨係賣家常嘢：\VUgras{茶葉}
\textsuperscript{(33)}、\VUgras{朱古力} \textsuperscript{(34)}、\VUgras{方糖} \textsuperscript{(37)}
同\VUgras{番梘} \textsuperscript{(38)}。呢間乜都有得賣：\VUgras{脆餅、}\VUgras{罐頭}
\textsuperscript{(35)}、\VUgras{洋酒} \textsuperscript{(36)}、\VUgras{冧酒、}\VUgras{拔蘭地、}
\VUgras{車厘酒、茴香酒}同\VUgras{橙皮酒。}嗰度仲有野味：\VUgras{野兔}
\textsuperscript{(39)}、\VUgras{畫眉} \textsuperscript{(40)}、\VUgras{鷓鴣} \textsuperscript{(41)}、
\VUgras{鵪鶉} \textsuperscript{(42)}、\VUgras{野鴨} \textsuperscript{(43)} 或者\VUgras{雉雞}
\textsuperscript{(44)}；外國生果一樣咁易搵，好似\VUgras{香蕉} \textsuperscript{(46)} 同
\VUgras{菠蘿。}

%%K t15-09-1 p
9. --- 雜貨舖個\VUgras{夥計} \textsuperscript{(45)} 好落力，你想要乜佢就攞乜：
\VUgras{果醬} \textsuperscript{(47)} 論斤定論份都得，\VUgras{豬油} \textsuperscript{(48)}、
\VUgras{酸青瓜} \textsuperscript{(49)}，或者鹹\VUgras{沙甸魚} \textsuperscript{(50)}。

%%K t15-09-2 p
噉樣對啲\VUgras{女客} \textsuperscript{(51)} 就好方便：家常嘢隔籬就係最罕有嘅時令貨
同最香甜嘅生果——多汁嘅\VUgras{梨} \textsuperscript{(52)}、\VUgras{李}
\textsuperscript{(53)}、\VUgras{提子} \textsuperscript{(54)}、\VUgras{桃} \textsuperscript{(55)}、
\VUgras{杏仁} \textsuperscript{(56)} 同\VUgras{合桃} \textsuperscript{(57)}。

%%K t15-tit-5 sub
\VUsaut{4.0mm}
\VUpk{10.2pt}{40}{客路。}
\VUsaut{3.0mm}

%%K t15-10-1 p
10. --- 啲\VUgras{米} \textsuperscript{(58)} 同啲\VUgras{扁豆} \textsuperscript{(59)} 就擺喺個
燻\VUgras{鯡魚} \textsuperscript{(60)} 箱隔籬；啲\VUgras{薯仔} \textsuperscript{(61)} 同啲
\VUgras{豆角} \textsuperscript{(63)} 就傍住啲\VUgras{蘋果} \textsuperscript{(62)}、\VUgras{無花果}
\textsuperscript{(64)}、\VUgras{梅乾} \textsuperscript{(65)} 同\VUgras{榛子} \textsuperscript{(66)}。
呢間舖仲有新鮮\VUgras{雞蛋} \textsuperscript{(67)} 同\VUgras{橄欖} \textsuperscript{(68)}；
所以啲\VUgras{籮} \textsuperscript{(70)} 同啲\VUgras{網袋} \textsuperscript{(73)} 一陣間就裝到滿瀉。

%%K t15-11-1 p
11. --- 呢間好舖頭啲\VUgras{咖啡} \textsuperscript{(72)}，喺啲\VUgras{女工人}
\textsuperscript{(69)} 同啲\VUgras{女廚}中間出咗名，而且係實至名歸——因為係個老
\VUgras{雜貨佬} \textsuperscript{(71)} 親自小心炒嘅。

%%K t15-tit-6 sub
\VUsaut{4.0mm}
\VUpk{10.2pt}{40}{睇熱鬧。}
\VUsaut{3.0mm}

%%K t15-12-1 p
12. --- 唔係個個嚟街市都係為買嘢。通去個大棚嗰條窄街好似幅畫噉熱鬧，
乜嘢人都見得到。一個好人好姐嘅\VUgras{屠房夥計} \textsuperscript{(75)} 推住架
\VUgras{手推車} \textsuperscript{(74)}；隔籬有兩個放緊假嘅兵，好威噉着住閃令令嘅制服：
一個\VUgras{驃騎兵} \textsuperscript{(76)}，藍色\VUgras{短上衣}配插羽嘅\VUgras{軍帽；}
另一個係\VUgras{祖阿夫下士，}衫袖鼓鼓哋，頭戴\VUgras{土耳其帽。}

%%K t15-13-1 p
13. --- \VUgras{麵包師傅} \textsuperscript{(80)} 企喺\VUgras{麵包舖} \textsuperscript{(78)}
門檻度，佢啱啱搓完\VUgras{麵包} \textsuperscript{(79)} 嘅麵糰，又由爐入面攞返出啲
\VUgras{牛角包} \textsuperscript{(81)}、啲\VUgras{圓包} \textsuperscript{(82)} 同啲\VUgras{大圓包}
\textsuperscript{(83)}——而家就擺喺櫥窗度。

%%K t15-14-1 p
14. --- 麵包舖樓上住住一位手工好好嘅\VUgras{車衫女} \textsuperscript{(84)}，佢請咗
幾個\VUgras{綉花女} \textsuperscript{(85)}。佢隔籬住住個\VUgras{洗熨婆} \textsuperscript{(86)}：
啲\VUgras{衫} \textsuperscript{(88)} 洗完晾乾之後，佢就上漿，再用個\VUgras{熨斗}
\textsuperscript{(87)} 熨平。

%%K t15-tit-7 sub
\VUsaut{4.0mm}
\VUpk{10.2pt}{40}{肉、魚同菜。}
\VUsaut{3.0mm}

%%K t15-15-1 p
15. --- 地下嗰個\VUgras{肉檔} \textsuperscript{(89)} 咩肉都有得賣：\VUgras{羊肉}
\textsuperscript{(90)}、\VUgras{牛肉} \textsuperscript{(94)} 或者\VUgras{牛仔肉} \textsuperscript{(92)}，
斬成\VUgras{羊髀、牛扒、燒肉}同\VUgras{肉排。}\VUgras{肉檔夥計} \textsuperscript{(93)}
就負責起晒啲肉，仲將啲\VUgras{骨髓骨} \textsuperscript{(96)} 另外擺開。肉檔門口有個
\VUgras{賣蠔婆} \textsuperscript{(97)} 賣\VUgras{蠔} \textsuperscript{(98)}，你想開就即場開畀你。

%%K t15-16-1 p
16. --- 呢啲商家全部要交牌費或者攤位費；所以個\VUgras{差人} \textsuperscript{(100)}
就毫不留情噉追住啲可憐嘅\VUgras{走鬼菜販} \textsuperscript{(99)} 趕——佢哋推住車仔賣
\VUgras{甘筍、水蘿蔔、蕪菁、韭蔥、}\VUgras{一紮紮蘆筍、新鮮青豆、菠菜、酸模、}
\VUgras{西洋菜}同\VUgras{芫茜，}同人哋搶生意。

%%K t15-tit-8 sub
\VUsaut{4.0mm}
\VUpk{10.2pt}{40}{睇住差人呀！}
\VUsaut{3.0mm}

%%K t15-17-1 p
17. --- 賣\VUgras{蒜頭} \textsuperscript{(101)} 同\VUgras{洋蔥}嗰個細路好清楚，佢一樣
唔准喺街市附近賣嘢。

%%K t15-17-1 p suite
佢一路走一路衝，差啲撞跌個\VUgras{賣龍蝦佬} \textsuperscript{(102)} 個籮——入面裝住
\VUgras{蟹、膏蟹}同\VUgras{蝦，}佢仲賣\VUgras{龍蝦}同\VUgras{大螯龍蝦。}

%%K t15-18-1 p
18. --- 個個都郁來郁去、嘈到拆天；但係咁都仲聽得清楚個沿街叫賣嘅
\VUgras{鑲玻璃佬} \textsuperscript{(103)} 把聲，同個\VUgras{賣炭佬} \textsuperscript{(104)} 嘅
叫喊——佢啱啱擺低架車仔一陣，去送一\VUgras{袋炭} \textsuperscript{(105)}。

\VUsaut{6.0mm}
\VUfilet{22mm}
